% !PRINTED 448
\weberpage{448}
Ist n\"amlich $x$ ein zweiz\"ahliger Pol, und geht $x$ durch $\omega$ in $x'$, durch $\Theta$ in $x''$ \"uber, so sind nicht nur $x'$ und $x''$ von $x$, sondern sie sind auch unter einander verschieden, weil, wenn sie identisch w\"aren, $x$ durch $\omega\Theta^3$ unge\"andert bliebe. Durch $\omega\Theta$ geht dann $x$ in eine vierte Lage $x'''$ \"uber, und $x'''$ ist sowohl von $x'$ als von $x''$ verschieden, weil sonst $x$ durch $\omega\Theta\omega = \Theta$ oder durch $\omega\Theta\Theta^{-1} = \omega$ unge\"andert bliebe. Da $x$ durch $\Theta^2$ unge\"andert bleibt, so geht $x$ durch $\Theta^3$ in $x''$, durch $\Theta^2\omega = \omega\Theta^2$ in $x'$, durch $\omega\Theta^3$ in $x'''$ \"uber.

Da wir auf jeder Hauptaxe zwei Punkte $c_1$, $c_2$ haben, so giebt es im Ganzen 42, und jeder von ihnen kann in jeden anderen durch Substitutionen der Gruppe \"ubergef\"uhrt werden. Daraus folgt, dass diese Pole nicht mehr als vierz\"ahlig sind. Daher der Satz:

\noindent 4.\ \gesperrt{Es giebt 21 achtz\"ahlige und 42 vierz\"ahlige Pole. Auf jeder Hauptaxe liegen vier achtz\"ahlige und zwei vierz\"ahlige Pole. Durch jeden achtz\"ahligen Pol gehen vier Hauptaxen. Jeder achtz\"ahlige Pol kann in jeden anderen achtz\"ahligen und jeder vierz\"ahlige Pol in jeden anderen vierz\"ahligen Pol \"ubergef\"uhrt werden.}

\begin{center}
\S.\,119.

\gesperrt{Die drei- und sechsz\"ahligen Pole.}
\end{center}

Wir wenden uns nun zur Betrachtung der Substitutionen dritter Ordnung und ihrer Pole, und w\"ahlen als Repr\"asentanten einer solchen Substitution

\[
\chi = \begin{pmatrix} 0 & 0 & 1\\ 1 & 0 & 0\\ 0 & 1 & 0 \end{pmatrix},
\]
deren Pole man aus

\[
\lambda\,x_1 = x_3, \quad \lambda\,x_2 = x_1, \quad \lambda\,x_3 = x_2
\]
erh\"alt. Es folgt daraus $\lambda^3 = 1$; also ist $\lambda$ eine dritte Einheitswurzel, und wenn daher $\varrho$ eine imagin\"are dritte Einheitswurzel bedeutet, so ergeben sich die drei Pole:

\weq{1}{\begin{array}{l}
x_1 = x_2 = x_3\\
x_1 = \varrho\,x_2 = \varrho^2 x_3\\
x_1 = \varrho^2 x_2 = \varrho\,x_3.
\end{array}}
